\documentclass[12pt]{article}
\usepackage[utf8]{inputenc}
\usepackage[T1]{fontenc}
\usepackage{hyperref}
\title{Shumcho Language Resources for K-12 Curriculum Development}
\author{Compiled by Claude Sonnet 5.5 (claude-sonnet-5-5) for Andr\'as Kornai}
\date{}
\begin{document}
\maketitle

\section{Introduction}
Shumcho (also Humcho; ISO 639-3 scu) is a West Himalayish (Kinnauric) Sino-Tibetan language of Kinnaur district, Himachal Pradesh, India. Sources report about 2,400 speakers (2022), against the roughly 3,200 in the survey prompt. It is spoken mainly in Poo tehsil, in villages such as Kanam, Labrang, Spillo, Shyaso, Taling and Rushkaling. It is distinct from neighbouring Sunam and Jangshung and is not mutually intelligible with Sunam. In some mixed villages it is spoken by Scheduled Caste groups. The only substantial documentation found is the academic work of Christian Huber. No school materials were found.

\section{Primary Educational Resources}
None found. No primer, textbook or school programme for Shumcho appeared in the searches.

\section{Linguistic and Reference Materials}
\begin{itemize}
\item Christian Huber's PhD dissertation (University of Vienna, 2017), \emph{Selected Topics in the Grammar of Shumcho/Humcho, Texts, and a Shumcho-English Glossary}. It is the principal source for grammar, texts and a glossary. It is listed in the publication list below; a full online copy was not confirmed. \url{https://www.oeaw.ac.at/fileadmin/Institute/PHA/PDF/Team/publikationen_huber_2026_04.pdf}
\item Huber's biography and CV at the Austrian Academy of Sciences (OeAW) Phonogrammarchiv. \url{https://www.oeaw.ac.at/fileadmin/Institute/PHA/PDF/Team/biographie_huber_2026_04.pdf}
\item Huber, ``Gender marking in Shumcho'', in a Benjamins volume (CILT 362). \url{https://benjamins.com/catalog/cilt.362.04hub}
\item OeAW project page, Researching and documenting endangered languages of Kinnaur. \url{https://www.oeaw.ac.at/en/phonogrammarchiv/research/endangered-languages-and-dialectology/researching-and-documenting-endangered-languages-of-kinnaur}
\item \emph{The Linguistic Landscape of the Indian Himalayas}, open access via OAPEN. \url{https://library.oapen.org/handle/20.500.12657/76653}
\item Wikipedia, Shumcho language. \url{https://en.wikipedia.org/wiki/Shumcho_language}
\end{itemize}
A descriptive grammar of Shumcho is listed as in preparation.

\section{Religious and Community Resources}
None found.

\section{Digital and Audio Resources}
The OeAW Phonogrammarchiv project above made field recordings. I did not confirm that any are publicly accessible online.

\section{Curriculum Development Context}
Shumcho has no known writing system or school use. Himachal Pradesh schools use Hindi and English. Any curriculum would be a community-based revitalization effort that builds on Huber's glossary and texts. It would need a community-approved orthography (probably Devanagari-based) and consultation with speakers, who in some villages belong to socially marginalised groups.
\end{document}